\documentclass[aps,prb,onecolumn,superscriptaddress,floatfix]{revtex4-2}

% 如果补充材料需要图表编号与主稿分开
\usepackage{chngcntr}
\counterwithin{figure}{section}
\counterwithin{table}{section}
\renewcommand{\thefigure}{S\arabic{figure}}
\renewcommand{\thetable}{S\arabic{table}}

\usepackage{graphicx}
\graphicspath{{figures/}}
\usepackage{amsmath,amssymb}

\begin{document}

\title{Supplemental Material for \\ ``A Neuroevolution Potential for Gallium Oxide: \\
Accurate and Efficient Modeling of Polymorphism and Swift Heavy-Ion Irradiation''}

\author{Yaohui Gu}
\affiliation{State Key Laboratory of Heavy Ion Science and Technology, Institute of Modern Physics, Chinese Academy of Sciences, Lanzhou 730000, China}
\affiliation{School of Nuclear Science and Technology, University of Chinese Academy of Sciences, Beijing 100049, China}
\author{Binbo Li}
\thanks{Y.G. and B.L. contributed equally to this work.}
\affiliation{State Key Laboratory of Heavy Ion Science and Technology, Institute of Modern Physics, Chinese Academy of Sciences, Lanzhou 730000, China}
\affiliation{School of Nuclear Science and Technology, University of Chinese Academy of Sciences, Beijing 100049, China}
\author{Linyang Jiang}
\affiliation{State Key Laboratory of Heavy Ion Science and Technology, Institute of Modern Physics, Chinese Academy of Sciences, Lanzhou 730000, China}
\affiliation{School of Nuclear Science and Technology, University of Chinese Academy of Sciences, Beijing 100049, China}
\author{Yuhui Hu}
\affiliation{State Key Laboratory of Heavy Ion Science and Technology, Institute of Modern Physics, Chinese Academy of Sciences, Lanzhou 730000, China}
\affiliation{School of Nuclear Science and Technology, University of Chinese Academy of Sciences, Beijing 100049, China}
\affiliation{School of Materials $\&$ Energy, Lanzhou University, Lanzhou, 730000, China}
\author{Wenqiang Liu}
\affiliation{State Key Laboratory of Heavy Ion Science and Technology, Institute of Modern Physics, Chinese Academy of Sciences, Lanzhou 730000, China}
\affiliation{School of Nuclear Science and Technology, University of Chinese Academy of Sciences, Beijing 100049, China}
\author{Lijun Xu}
\affiliation{State Key Laboratory of Heavy Ion Science and Technology, Institute of Modern Physics, Chinese Academy of Sciences, Lanzhou 730000, China}
\affiliation{School of Nuclear Science and Technology, University of Chinese Academy of Sciences, Beijing 100049, China}
\author{Pengfei Zhai}
\affiliation{State Key Laboratory of Heavy Ion Science and Technology, Institute of Modern Physics, Chinese Academy of Sciences, Lanzhou 730000, China}
\affiliation{School of Nuclear Science and Technology, University of Chinese Academy of Sciences, Beijing 100049, China}
\author{Jie Liu}
\affiliation{State Key Laboratory of Heavy Ion Science and Technology, Institute of Modern Physics, Chinese Academy of Sciences, Lanzhou 730000, China}
\affiliation{School of Nuclear Science and Technology, University of Chinese Academy of Sciences, Beijing 100049, China}
\author{Jinglai Duan}
\thanks{Corresponding author. Email: j.duan@impcas.ac.cn}
\affiliation{State Key Laboratory of Heavy Ion Science and Technology, Institute of Modern Physics, Chinese Academy of Sciences, Lanzhou 730000, China}
\affiliation{School of Nuclear Science and Technology, University of Chinese Academy of Sciences, Beijing 100049, China}
\affiliation{Advanced Energy Science and Technology Guangdong Laboratory, Huizhou 516000, China}

\date{\today}

\maketitle

\section{Extra Figures}

\begin{figure}[ht]
\centering
\includegraphics[width=0.8\linewidth]{S1.png}
\caption{Energy-volume curves for the $\alpha$, $\beta$ and $\gamma$ phases 
predicted by zero-shot NEP (solid circles), NEP with augmented training set (solid triangles), 
and tabGAP (solid squares), compared with DFT reference 
data (open circles). Different phases are color-coded as indicated.}
\label{fig:S1}
\end{figure}

\begin{figure}[ht]
\centering
\includegraphics[width=0.8\linewidth]{S2.png}
\caption{\textbf{(a)} Longitudinal cross-sectional view of the SHI-irradiated $\beta$-Ga$_2$O$_3$ 
atomic structure with $S_e = 19.18$~keV/nm, the vertical direction is along [201] direction. 
\textbf{(b)} Simulated diffraction pattern for the pristine $\beta$ phase.
\textbf{(c)} Simulated diffraction pattern for the $\gamma$ phase.}
\label{fig:S2}
\end{figure}

\begin{figure}[ht]
\centering
\includegraphics[width=0.75\linewidth]{figureS3_transverse_part1.png}
\caption{Transverse cross-sectional views of ion tracks in $\beta$-Ga$_2$O$_3$ 
for electronic energy losses ranging from $S_e = 11.29$~keV/nm to 
$S_e = 26.05$~keV/nm. The measured track diameters are indicated for each case.}
\label{fig:S3}
\end{figure}

\begin{figure}[ht]
    \centering
    \includegraphics[width=0.75\linewidth]{figureS3_transverse_part2.png}
    \caption{Transverse cross-sectional views of ion tracks in $\beta$-Ga$_2$O$_3$ 
    for electronic energy losses ranging from $S_e = 30.08$~keV/nm to 
    $S_e = 42.92$~keV/nm. The measured track diameters are indicated for each case.}
    \label{fig:S4}
\end{figure}

\end{document}
